\subsection*{Underwater Object Detection}
Underwater object detection has received increasing attention because of its relevance to
marine robotics and ecological monitoring. Recent methods mainly extend one-stage detectors
with improved multiscale fusion, lightweight backbones, or attention mechanisms to handle
complex underwater backgrounds and real-time constraints.

\subsection*{Underwater Image Enhancement for Detection}
Image enhancement is widely used to mitigate underwater degradation before downstream visual
tasks. Nevertheless, enhancement does not uniformly improve detector behavior. Its benefits
depend on degradation type, target scale, and the reliability of the recovered structures.
This motivates adaptive cooperation between raw and enhanced domains instead of relying on a
single fixed preprocessing strategy.

\subsection*{Small Object Detection}
Small object detection remains difficult because small targets occupy limited pixels and are
easily overwhelmed by background noise. Common strategies include feature pyramid redesign,
high-resolution detail preservation, shallow feature injection, and task-specific attention.

\subsection*{Multi-Branch Feature Interaction}
Multi-branch architectures are often used to encode heterogeneous cues. In underwater
detection, raw and enhanced images provide complementary information. The key challenge is
not only how to fuse them, but how to prevent noisy or unstable cross-domain interactions.
This paper addresses that issue through cross-residual exchange and quality-aware fusion.
